\documentclass[RNAAS]{aastex631}

% Numbers are \input from the pipeline-generated macros -- none typed by hand.
% Auto-generated by jansky_research.ucdv._write_macros -- do not edit.
\newcommand{\ucdvSource}{synthetic recover-a-known + real RACS-low2/low3 forced-V census}
\newcommand{\ucdvSynNTargets}{2000}
\newcommand{\ucdvSynNPlanted}{6}
\newcommand{\ucdvSynNRecovered}{6}
\newcommand{\ucdvSynNFalse}{0}
\newcommand{\ucdvRealNReyle}{14915}
\newcommand{\ucdvRealNBoth}{12751}
\newcommand{\ucdvRealDmax}{107.8}
\newcommand{\ucdvRealNCensus}{9805}
\newcommand{\ucdvRealFullHours}{156.0}
\newcommand{\ucdvRealSzeroOk}{80}
\newcommand{\ucdvRealSzeroN}{80}
\newcommand{\ucdvRealSzeroSec}{22.0}
\newcommand{\ucdvRealConeN}{400}
\newcommand{\ucdvRealConePos}{0.51}
\newcommand{\ucdvRealConeStd}{0.96}
\newcommand{\ucdvRealCtwoN}{480}
\newcommand{\ucdvRealCtwoComp}{0.965}
\newcommand{\ucdvRealCtwoRatio}{0.982}
\newcommand{\ucdvRealCthreeN}{67}
\newcommand{\ucdvRealCthreeRatio}{1.021}
\newcommand{\ucdvRealCthreeFlips}{0}
\newcommand{\ucdvRealNTargets}{9805}
\newcommand{\ucdvRealNSlots}{19600}
\newcommand{\ucdvRealNSane}{19576}
\newcommand{\ucdvRealNCand}{0}
\newcommand{\ucdvRealNHeld}{1}
\newcommand{\ucdvRealNPairs}{9771}
\newcommand{\ucdvRealNBeyondSpec}{0}
\newcommand{\ucdvRealSigmaV}{0.151}
\newcommand{\ucdvRealLimit}{0.57}
\newcommand{\ucdvRealNullPos}{0.501}
\newcommand{\ucdvRealNullStd}{1.009}
\newcommand{\ucdvRealNeffFourteen}{2}
\newcommand{\ucdvRealFracFourteen}{1.49787}
\newcommand{\ucdvRealNeffFifteen}{860}
\newcommand{\ucdvRealFracFifteen}{0.00348}
\newcommand{\ucdvRealNeffSixteen}{19103}
\newcommand{\ucdvRealFracSixteen}{0.00016}


\newcommand{\code}[1]{\texttt{#1}}

\begin{document}

% DRAFT (2026-10-17). The title carries the verdict and waits for the census; the results
% paragraph and the abstract's result sentences are written only after the census is analysed.
\title{Circular Polarization of Gaia Ultracool Dwarfs in Two RACS Epochs at 0.9~GHz}

\author[0009-0008-3289-4447]{Joseph Barbere}
\affiliation{Independent researcher}

\begin{abstract}
We measured Stokes~$V$ at the positions of \ucdvRealNCensus\ ultracool dwarfs from the
\citet{reyle2018} Gaia DR2 sample within \ucdvRealDmax~pc, in the RACS-low2 (887.5~MHz) and RACS-low3
(943.5~MHz) epochs of the Rapid ASKAP Continuum Survey. % RESULT SENTENCES PENDING THE CENSUS.
\end{abstract}

\section{Sample and measurement}
\label{sec:method}
Ultracool-dwarf radio emission is coherent and often strongly circularly polarized, so a
Stokes~$V$ survey selects it against the unpolarized extragalactic background. The blind
VLASS search of the \citet{reyle2018} candidates at 3~GHz used Stokes~$I$ alone
\citep{perger2025}, and the RACS-low2 circular-polarization catalogue measures $V$ only at
Stokes~$I$ catalogue positions \citep{duchesne2026}. We instead measure $V$ at every dwarf.

The parent sample is the \citet{reyle2018} table of Gaia DR2 \citep{gaiadr2} candidates.
Of its \ucdvRealNReyle\ entries, \ucdvRealNBoth\ fall inside both a low2 and a low3 Stokes~$V$ image on CASDA; RACS-low1
has no $V$ products there. The fetch rate measured on \ucdvRealSzeroN\ test pairs
(\ucdvRealSzeroOk\ succeeded, \ucdvRealSzeroSec~s per pair) projected \ucdvRealFullHours~h
for the full sample; a rule fixed beforehand then set a distance limit of \ucdvRealDmax~pc,
chosen from parallax alone. For each dwarf and epoch we
retrieved a $0.06^\circ$ $I$ and $V$ cutout from one observation \citep{mcconnell2020}
with \code{astroquery} \citep{astroquery}, moved the Gaia position to the observation date
with its proper motion, and read $I$ and $V$ at the pixel containing it. The noise is the
standard deviation in an annulus of 30--90\arcsec.

A target-epoch counts as a candidate if $|V|/\sigma_V\geq5$, if $|V|/I$ exceeds the local
leakage floor when Stokes~$I$ is itself detected, and if no Stokes~$I$ source brighter than
100~mJy lies within 3\arcmin. The upper limit on $|V|$ is $|V|+3\sigma_V$.

\section{Controls}
\label{sec:controls}
The thresholds and pass criteria were fixed before any target value was examined. The checks
used the \ucdvRealSzeroN\ test pairs, whose dwarfs come from the half of the sample (split by a
hash of the Gaia identifier) that sets no threshold. Forced photometry at
\ucdvRealConeN\ blank positions 2\arcmin\ from the targets gave positive $V$ in a fraction
\ucdvRealConePos\ of cases with $\mathrm{std}(V/\sigma_V)=\ucdvRealConeStd$, as expected
for noise; a peak search instead of a fixed pixel would bias this fraction towards one. Point
sources injected into the real cutouts (\ucdvRealCtwoN\ per level) were recovered at a rate of
\ucdvRealCtwoComp\ at $7\sigma_V$, with a median recovered-to-injected flux ratio of
\ucdvRealCtwoRatio\ above $10\sigma_V$. For the \ucdvRealCthreeN\ stars of the RACS-low2
circular-polarization catalogue \citep{duchesne2026}, measured in their own observations, the
median ratio of our $V$ to the catalogue value is \ucdvRealCthreeRatio, with
\ucdvRealCthreeFlips\ sign reversals.

\section{Results}
\label{sec:results}
% PENDING THE CENSUS: candidates and vetting, the limit table, the fraction of snapshots above
% L_nu (divided by the snapshots deep enough to see it), the two-epoch comparison, and WISE
% J0623-0456 \citep{rose2023} reported as a non-gating case.

\software{\code{jansky-research} \citep{janskyresearch}, \code{jansky} \citep{jansky},
\code{astroquery} \citep{astroquery}, NumPy, SciPy, Astropy. Analysis assisted by Claude
\citep{anthropicclaude}. The scientific kernels are covered by an offline test suite; the
network and real-data drivers are not unit-tested.}

\bibliography{refs}
\bibliographystyle{aasjournal}

\end{document}
